\documentclass[11pt]{article}
\usepackage{amsmath,amssymb,amsthm}
\newtheorem{theorem}{Theorem}
\newtheorem{lemma}{Lemma}
\begin{document}
\title{Superlinear Readouts of Permutation Sheets with Free Assignment Costs}
\author{Independent gluing audit}
\date{November 23, 2026}
\maketitle

\paragraph{Scope and verification status.}
The theorem below is an unconditional, complete written statement about
\emph{free linear assignment costs}. It proves that the factorial sheet
storage is sufficient for superlinear sign decoding in this larger query
class. It does not represent these arbitrary costs by admissible reciprocal
queries. Consequently it supplies neither the missing local compiler nor
a new bound on head complexity. The proof is not formalized in Lean.

\begin{theorem}
For every integer $H\geq64$, there exist
\[
d\geq\frac{H\log H}{80}
\]
real matrices $C^1,\ldots,C^d\in\mathbb R^{H\times H}$ and thresholds
$t_1,\ldots,t_d$ such that for every sign vector
$s\in\{-1,1\}^d$ there is a permutation $\pi\in\mathfrak S_H$ with
\[
s_i\left(\sum_{h=1}^H C^i_{h,\pi(h)}-t_i\right)>0
\quad\text{for all }i=1,\ldots,d.
\]
The matrices and thresholds are fixed independently of the sign vector.
\end{theorem}

\begin{proof}
We give the packing, separation, and extraction arguments explicitly.

\paragraph{Permutation packing.}
Let $r=\lfloor H/4\rfloor$. A Hamming ball of radius $r$ in
$\mathfrak S_H$ has at most
\[
\sum_{k=0}^{r}\binom Hk k!
\leq(H+1)H^{H/4}
\]
elements: after selecting the $k$ changed positions there are at most
$k!$ ways to permute their values. Greedily choose a permutation and
delete its radius-$r$ ball until none remain. The resulting code
$\mathcal P$ has pairwise Hamming distance greater than $r$, hence
greater than $H/4$, and cardinality $N$ satisfying
\[
N\geq\frac{H!}{(H+1)H^{H/4}},\qquad
\log N\geq\frac H4\log H-\frac H2\log2-\log(H+1)
\geq\frac H8\log H.
\]
For the factorial estimate, at least $H/2$ factors of $H!$ are at least
$H/2$. For the final inequality put
$g(x)=x\log x/8-x\log2/2-\log(x+1)$.
At $x=64$ one has $g(64)=\log(2^{16}/65)>0$; for $x\geq64$,
\[
g'(x)\geq\frac{6\log2+1}{8}-\frac{\log2}{2}-\frac1{65}>0.
\]
Thus the displayed inequality holds for every stated $H$.

\paragraph{Geometric separation.}
Let $P_\pi$ be the permutation matrix of $\pi$, let $J$ be the all-one
matrix, and put $v_\pi=P_\pi-J/H$. Frobenius inner products give
\[
\|v_\pi\|_F^2=H-1,\qquad
\frac{\langle v_\pi,v_\tau\rangle_F}{H-1}
=1-\frac{\operatorname{dist}(\pi,\tau)}{H-1}<\frac34
\quad(\pi\ne\tau\text{ in }\mathcal P).
\]
Their pairwise angle is therefore greater than $\pi/6$.
If $G$ is a matrix of independent standard Gaussian entries, rotational
invariance in the two-dimensional span of $v_\pi,v_\tau$ shows
\[
\Pr\bigl(\operatorname{sgn}\langle G,v_\pi\rangle_F
\ne\operatorname{sgn}\langle G,v_\tau\rangle_F\bigr)>\frac16.
\]
Indeed the two opposite separating angular wedges have total angle
$2\theta$ in a circle of angle $2\pi$, where $\theta$ is their angle.
Each inner product is nonzero with probability one.

Take $q=\lceil18\log N\rceil$ independent Gaussian matrices.
The probability that any pair of codewords has the same $q$ signs is
at most
\[
N^2(5/6)^q\leq N^2e^{-q/6}\leq N^{-1}<1.
\]
There is consequently a deterministic choice of these $q$ matrices
whose sign vectors distinguish all $N$ codewords and contain no zero
values. Each sign is a free assignment-cost query: use $C^i=G_i$ and
$t_i=\langle G_i,J/H\rangle_F$.

\paragraph{Extraction of independent bits.}
Let $\mathcal A\subseteq\{-1,1\}^q$ be the resulting $N$ distinct
sign vectors. If its largest shattered coordinate subset has size $d$,
the Sauer counting inequality gives
\[
N\leq\sum_{j=0}^d\binom qj\leq(eq/d)^d\qquad(d\geq1).
\]
For completeness, the first inequality follows by induction on the
number of coordinates: partition according to the last coordinate;
the union of the projected families has shattering dimension at most
$d$, and their intersection at most $d-1$. Pascal's identity completes
the induction. The second inequality follows by multiplying the sum by
$(d/q)^d$, bounding each term by its counterpart in
$(1+d/q)^q$, and using $(1+d/q)^q\leq e^d$.

Since $\log N\geq1$, one has $q\leq19\log N$.
If $d<\log N/10$, monotonicity of
$x\mapsto x\log(eq/x)$ for $0<x\leq q$ would imply
\[
\log N\leq d\log(eq/d)
\leq\frac{\log N}{10}\log(190e)<\log N,
\]
a contradiction. Here $\log(190e)<10$ follows from $190<2^9<e^9$.
Thus $d\geq\log N/10\geq H\log H/80$.
Choose a shattered coordinate subset and retain just those matrices
and thresholds. This proves the theorem with strict signs.
\end{proof}

\paragraph{Connection to regular sheets.}
At the regular reciprocal sheets of
\texttt{reverse-joint-full-sheets.tex}, the slot parameters satisfy
$(u_h,a_h)=(\lambda_{\pi(h)},\mu_{\pi(h)})$.
A permitted extra reciprocal query has sheet values
\[
\sum_h\frac{A_h+\mu_{\pi(h)}}{d_h-\lambda_{\pi(h)}}+c.
\]
These are free assignment costs only with the restricted matrices
$C_{hj}=(A_h+\mu_j)/(d_h-\lambda_j)$.
The theorem above permits arbitrary $C_{hj}$ and does not establish
that this restricted family contains its separating queries.
For any genuinely permitted strict readouts, regularity would let the
original $2H$ continuous query signs vary in a sufficiently small
neighborhood while preserving all extra signs. The unresolved transfer
is therefore quantitative separation by the restricted admissible
query family, not the existence of factorial storage or the final
inverse-function step.

\end{document}
